% Конспект изменений: регуляторы лидер/преследователь, визуализатор без V-Bots
% Кодировка: UTF-8 (pdflatex + inputenc или xelatex/lualatex)

\documentclass[12pt,a4paper]{article}
\usepackage[utf8]{inputenc}
\usepackage[T2A]{fontenc}
\usepackage[russian]{babel}
\usepackage{amsmath,amssymb}
\usepackage{enumitem}
\usepackage{hyperref}
\usepackage{geometry}
\geometry{margin=2.5cm}

\title{Конспект изменений в проекте симуляции ансамбля БПЛА}
\author{Лидер--преследователь, регуляторы, визуализатор}
\date{}

\begin{document}
\maketitle

\section{Обзор}

В проекте реализован сценарий «лидер -- преследователь»: лидер движется вперёд и назад по заданной программе, преследователь удерживает заданное смещение относительно лидера с помощью PID-регуляторов по каналам roll и pitch. Ниже перечислены все внесённые изменения: доработка регуляторов (пропорциональная и дифференциальная составляющие), настройка контура управления, добавление визуализатора, работающего без 3D-симуляции V-Bots (Webots).

\section{Регуляторы}

\subsection{Лидер}

Лидер \emph{не} использует регулятор по позиции. Управление задаётся программно:
\begin{itemize}[noitemsep]
  \item Режим полёта: \texttt{POS\_HOLD} (удержание позиции с интерпретацией стиков как углов крена/тангажа).
  \item Движение вперёд: постоянная команда по каналу pitch (\texttt{pitch = 1400}, наклон вперёд от нейтрали 1500) в течение заданного времени (\texttt{leader-duration}, по умолчанию 10 с).
  \item Остановка: переход в режим \texttt{BRAKE}, пауза, возврат в \texttt{POS\_HOLD} с нейтральными стиками.
\end{itemize}
Таким образом, у лидера нет пропорциональной или дифференциальной составляющей по ошибке позиции --- только программное задание «вперёд / стоп».

\subsection{Преследователь: PID по roll и pitch}

Преследователь использует два независимых PID-регулятора: по ошибке по оси X (North) выдаётся управление по pitch, по ошибке по оси Y (East) --- по roll. Целевая позиция --- позиция лидера с заданным смещением (например, 2 м по Y).

\subsubsection{Структура регулятора}

Класс \texttt{PIDRegulator} реализует закон:
\begin{equation}
  u = K_p\, e + K_i \int e\, dt + K_d\, \frac{de}{dt},
\end{equation}
с ограничениями: интеграл ограничен (\texttt{integral\_limit}), выход ограничен (\texttt{output\_limit}). По умолчанию для преследователя: \texttt{integral\_limit = 100}, \texttt{output\_limit = 200} (по каналу pitch/roll в единицах RC 1300--1700).

\subsubsection{Пропорциональная составляющая (P)}

\begin{itemize}[noitemsep]
  \item \textbf{Назначение:} реакция на текущую ошибку позиции; при \texttt{error\_x > 0} (лидер впереди) преследователь должен давать наклон вперёд (pitch $<$ 1500).
  \item \textbf{Изменение:} коэффициент $K_p$ увеличен с 1 до \textbf{8} по умолчанию. При $K_p = 1$ при большой ошибке (например 50 м) выход P составлял 50, что давало слабый наклон и отставание преследователя по X. При $K_p = 8$ преследователь быстрее выходит на «полный газ вперёд» и догоняет лидера.
  \item \textbf{Параметр CLI:} \texttt{--kp} (по умолчанию 8.0).
\end{itemize}

\subsubsection{Интегральная составляющая (I)}

\begin{itemize}[noitemsep]
  \item По умолчанию \texttt{--ki 0}; при необходимости можно включить для устранения статической ошибки.
  \item Используется anti-windup: интеграл ограничен \texttt{integral\_limit}.
\end{itemize}

\subsubsection{Дифференциальная составляющая (D)}

\begin{itemize}[noitemsep]
  \item \textbf{Назначение:} демпфирование при сближении с целью; уменьшение перелётов (overshoot) и осцилляций.
  \item \textbf{Изменение:} введена дифференциальная составляющая с коэффициентом по умолчанию \textbf{$K_d = 10$} (\texttt{--kd 10}).
  \item \textbf{Сглаживание D:} чтобы производная по ошибке не давала всплесков из-за шума, введён экспоненциальный фильтр: \texttt{derivative\_alpha = 0.7} (настраиваемый при создании регулятора). Фильтр: $d_{filt} \leftarrow \alpha\, d_{filt} + (1-\alpha)\, \frac{de}{dt}$.
  \item \textbf{Параметр CLI:} \texttt{--kd} (по умолчанию 10.0).
\end{itemize}

\subsubsection{Фиксированный шаг времени (dt) и частота контура}

\begin{itemize}[noitemsep]
  \item При включённом ограничении частоты (\texttt{--control-hz 50}) в \texttt{update()} PID передаётся \textbf{фиксированный} $dt = 0.02$ с, чтобы интеграл и производная считались стабильно, без джиттера.
  \item Частота цикла преследования по умолчанию --- \textbf{50 Гц} (\texttt{--control-hz 50}); при \texttt{--control-hz 0} ограничение снимается.
\end{itemize}

\subsection{Сводка параметров по умолчанию (преследователь)}

\begin{center}
\begin{tabular}{|l|c|l|}
  \hline
  Параметр & Значение & Описание \\
  \hline
  \texttt{--kp} & 8.0 & P-коэффициент \\
  \texttt{--ki} & 0.0 & I-коэффициент \\
  \texttt{--kd} & 10.0 & D-коэффициент (демпфирование) \\
  \texttt{--control-hz} & 50.0 & Частота контура (Гц); 0 = без ограничения \\
  \hline
\end{tabular}
\end{center}

\section{Данные и контур обратной связи}

\begin{itemize}[noitemsep]
  \item \textbf{Запрос частой выдачи позиции:} при инициализации каждого дрона отправляется \texttt{SET\_MESSAGE\_INTERVAL} для \texttt{LOCAL\_POSITION\_NED} с интервалом 20\,000 мкс (50 Гц), чтобы автопилот отдавал позицию с предсказуемой частотой, а не 1--4 Гц по умолчанию.
  \item \textbf{Расширенное логирование:} в CSV пишутся $t$, \texttt{leader\_x}, \texttt{leader\_y}, \texttt{follower\_x}, \texttt{follower\_y}, \texttt{error\_x}, \texttt{error\_y}, \texttt{follower\_vx}, \texttt{follower\_vy}, \texttt{loop\_dt} для последующего анализа и построения графиков.
\end{itemize}

\section{Визуализатор без V-Bots}

Добавлена возможность наблюдать траектории лидера и преследователя в плоскости X--Y \emph{без} запуска 3D-симуляции Webots (V-Bots).

\subsection{Визуализация в реальном времени по логу}

Скрипт \texttt{plot\_realtime.py} читает лог-файл (\texttt{logs/two\_drones\_log.csv}), который в реальном времени пишет сценарий \texttt{leader\_forward\_one\_move.py}, и обновляет график позиций в плоскости X--Y.

\begin{itemize}[noitemsep]
  \item \textbf{Режим «в файл» (без окна):} \texttt{--output FILE} --- кадр сохраняется в изображение (например \texttt{logs/realtime\_xy.png}) с заданным интервалом (по умолчанию 0.15 с). Не требует Qt/tkinter; подходит для окружений, где GUI недоступен (например Wayland без плагина xcb). Для автообновления картинки в просмотрщике: \texttt{feh -R 0.2 logs/realtime\_xy.png}.
  \item \textbf{Режим с окном:} при запуске без \texttt{--output} скрипт пытается использовать GUI-бэкенд (Qt5Agg, TkAgg и др.); при отсутствии или ошибке рекомендуется перейти на режим \texttt{--output}.
\end{itemize}

Пример:
\begin{verbatim}
# Терминал 1: сценарий
python leader_forward_one_move.py

# Терминал 2: визуализатор в файл (без V-Bots)
python plot_realtime.py logs/two_drones_log.csv --output logs/realtime_xy.png
feh -R 0.2 logs/realtime_xy.png
\end{verbatim}

\subsection{Построение графиков по готовому логу}

Скрипт \texttt{plot\_log\_coords.py} строит по CSV четыре графика: X и Y от времени (лидер и преследователь), ошибки по осям от времени, траектория в плоскости XY. При указании \texttt{-o plot.png} рисунок сохраняется в файл без открытия окна (бэкенд Agg), что также позволяет обходиться без GUI и без V-Bots.

\subsection{Единый лаунчер}

Лаунчер (\texttt{launch\_simulation.py} или \texttt{launch.py}) позволяет из одного места выбрать: симуляцию с Webots или только SITL, сценарий (лидер вперёд, квадрат и др.) и включение визуализатора. При выборе «только SITL» + визуализатор траектории отображаются по логу, без запуска V-Bots.

\section{Краткая сводка изменений}

\begin{enumerate}
  \item \textbf{Лидер:} программное задание (постоянный pitch вперёд, BRAKE, возврат); регуляторов по позиции нет.
  \item \textbf{Преследователь:} введены и настроены пропорциональная (P) и дифференциальная (D) составляющие PID; $K_p$ увеличен до 8, $K_d = 10$ по умолчанию; сглаживание D-члена (derivative\_alpha); фиксированный $dt$ при частоте контура 50 Гц.
  \item \textbf{Данные:} запрос выдачи \texttt{LOCAL\_POSITION\_NED} с частотой 50 Гц; расширенное логирование в CSV.
  \item \textbf{Визуализатор без V-Bots:} \texttt{plot\_realtime.py} с режимом \texttt{--output} для сохранения кадра в файл; \texttt{plot\_log\_coords.py -o} для графиков по логу без окна; лаунчер с опцией «только SITL» + визуализатор.
\end{enumerate}

\end{document}
